% Generated by paper/build.py; numerical source remains unchanged.
% Generated by scripts/summarize_study.py; do not transcribe historical metrics here.
\subsection{Degree-one persistence and scale controls}
Degree-one features and all matched comparators completed on the full audited cohort. Table~\ref{tab:persistence} compares each addition with the graph-plus-degree-zero model using the same sheaf construction. Lower-endpoint, upper-endpoint, and persistent additions each supply twelve spectral summaries, so their feature counts match.


\begin{table}[htbp]\centering\small
\caption{Matched degree-one additions under sequence-controlled folds. The changes are paired mean per-protein PCC differences from the corresponding graph-plus-degree-zero model.}
\label{tab:persistence}
\begin{tabular}{llrr}\toprule
Sheaf & Degree-one addition & Mean PCC & Paired change [95\% CI]\\\midrule
Identity & Ordinary lower & 0.630 & +0.0002 [-0.0011,\allowbreak +0.0015] \\
Identity & Ordinary upper & 0.629 & -0.0002 [-0.0015,\allowbreak +0.0012] \\
Identity & Persistent & 0.629 & -0.0002 [-0.0015,\allowbreak +0.0011] \\
Geometric & Ordinary lower & 0.630 & -0.0006 [-0.0024,\allowbreak +0.0010] \\
Geometric & Ordinary upper & 0.629 & -0.0017 [-0.0034, -0.0002] \\
Geometric & Persistent & 0.630 & -0.0006 [-0.0022,\allowbreak +0.0010] \\
Center-zero & Ordinary lower & 0.627 & -0.0004 [-0.0022,\allowbreak +0.0013] \\
Center-zero & Ordinary upper & 0.628 & +0.0005 [-0.0010,\allowbreak +0.0021] \\
Center-zero & Persistent & 0.627 & +0.0000 [-0.0018,\allowbreak +0.0017] \\
\bottomrule\end{tabular}\end{table}


\begin{samepage}
The persistent-minus-upper-ordinary comparisons are $+0.0000$ (95\% CI $[-0.0013,\allowbreak +0.0012]$) for identity; $+0.0011$ (95\% CI $[-0.0005,\allowbreak +0.0029]$) for geometric; $-0.0005$ (95\% CI $[-0.0024,\allowbreak +0.0012]$) for center-zero. These comparisons isolate the value of the constrained two-scale operator from the simpler alternative of accessing the larger filtration radius. The point estimates describe predictive effects of the summarized spectra; the operators can differ algebraically even when their predictive performances are close.
\par\end{samepage}

The descriptor audit confirms that persistence changes the represented information. Identity and all-one geometric restrictions have identical numerical nullities at every tested neighborhood and scale, as expected from the invertible gauge. For the identity sheaf, persistent degree-one nullity is nonzero in 28.58\% of neighborhoods for $(3,4)$ and 0.361\% for $(5,6)$. These descriptive frequencies weight all retained residue-centered neighborhoods equally; they are not protein-level prediction statistics. The detailed comparison tables below report both endpoint counts and center-zero counts.

Across the three sheaves, two scale pairs, and two endpoint comparisons, the five positive-spectrum summary vectors differ in 98.94--99.98\% of neighborhoods (relative tolerance $10^{-8}$, absolute tolerance $10^{-10}$). This compares stored summaries, not full matrices of potentially different dimensions; equal summary vectors would not establish spectral equivalence.

For the graph-plus-center-zero degree-zero model, retaining only nullities gives mean PCC 0.627, while retaining only positive-spectrum summaries gives 0.627. Single-scale models at radii 3, 4, 5, and 6 \AA\ give 0.628, 0.625, 0.628, 0.626, respectively. These fixed ablations retain the graph features; they do not attribute a causal effect to topology or identify an optimal radius. The complete intervals and paired changes are included in the detailed comparison tables below.

The initial degree-zero pass used 33.5 summed worker minutes; the degree-one addition used 70.2, and the added upper-endpoint controls used 55.6. Four feature workers were used. The initial pass preceded the modified-residue amendment, so the degree-one timing also includes recalculation of degree-zero features for the 24 affected structures. These measured costs include matrix construction and spectral summarization, and are distinct from model-fitting time.

